\documentclass[12pt,a4paper]{article}
\usepackage[utf8]{inputenc}
\usepackage[T1]{fontenc}
\usepackage{amsmath,amssymb}
\usepackage{booktabs}
\usepackage{longtable}
\usepackage{geometry}
\geometry{margin=2.5cm}
\usepackage{enumitem}
\usepackage{hyperref}
\usepackage{xcolor}
\usepackage{caption}

\title{Classification Protocol: Mapping Fiscal Policy Measures\\to Transmission Channels (DI/CP/H)}
\author{Working Document --- Aulis Pesenti}
\date{February 2026}

\begin{document}
\maketitle

\section{Classification Principle}

The classification assigns each fiscal policy measure in the dataset to one of three
transmission channel categories based on the theoretical framework developed in
Section~2 of the paper. The categories correspond to structurally distinct channels
through which fiscal interventions affect the state variables of the planner's problem:

\begin{itemize}[leftmargin=2em]
  \item \textbf{Demand Injection (DI):} Measures that directly increase the flow of
    disposable income to households or aggregate expenditure in the economy. These
    transmit through the multiplier channel $\alpha_F^{DI}$ in Equation~(14), shifting
    the output level by a constant amount per unit deployed.
  \item \textbf{Capacity Preservation (CP):} Measures that preserve productive
    structures---firm--worker matches, firm solvency, supply chain integrity---thereby
    reducing the persistence of the output gap. These transmit through the parameter
    $\eta$ in Equation~(14), counteracting lockdown-induced hysteresis ($\psi S_k$).
  \item \textbf{Health (H):} Pandemic-specific health expenditures that are not a
    policy instrument in the planner's fiscal composition problem but an endogenous
    cost driven by the epidemiological state $\theta_k$. These enter the model through
    $c_H \theta_k$ in the debt equation~(16).
\end{itemize}

\noindent The classification operates at the \texttt{PolicyCode} level (49 codes as defined
in the Codebook, Table~2). This ensures reproducibility and eliminates subjective
judgment at the individual measure level, with one exception: the residual category
\texttt{other} (41 observations, all policy extensions with zero fiscal weight) is
classified case-by-case based on the measure description.

\section{Decision Rules}

The following hierarchy governs classification decisions when a measure's transmission
channel is not immediately apparent from its \texttt{PolicyCode} definition:

\begin{enumerate}[leftmargin=2em]
  \item \textbf{Primary criterion: transmission mechanism.} The decisive question
    is \emph{how} the measure affects the economy, not \emph{who} receives the funds.
    A wage subsidy paid to firms that maintains employment relationships is CP even
    though it ultimately supports worker income. A direct cash transfer to households
    is DI even if the household uses it to pay rent to a business.
  \item \textbf{Recipient as tiebreaker.} When the transmission mechanism is ambiguous,
    the recipient disambiguates: measures targeting firms, employers, or productive
    structures default to CP; measures targeting households, individuals, or consumers
    default to DI.
  \item \textbf{Tax deferrals.} All tax deferrals (IMF Category~3) are classified as
    CP regardless of whether the recipient is a business or individual. The
    rationale: deferrals operate as zero-interest short-term government loans that
    preserve liquidity and prevent insolvency cascades. They generate no permanent
    change in disposable income (the deferred amount is repaid) and thus have no
    sustained multiplier effect through $\alpha_F^{DI}$. Their primary transmission
    is capacity preservation---preventing liquidity-constrained agents from defaulting
    on obligations during the acute crisis phase.
  \item \textbf{Health expenditures.} All measures classified under PolicyCategory~5
    (Health expenditure measures) and 5.5 (Revenue measures to promote availability
    of medical items) are assigned to H, regardless of whether they involve
    expenditure or revenue instruments.
  \item \textbf{Code-level homogeneity.} Classification is applied uniformly within
    each \texttt{PolicyCode}. For codes with heterogeneous content (notably Code~39),
    the dominant transmission mechanism determines the classification, and an
    ambiguity flag is set for robustness checks.
\end{enumerate}


\section{Complete Classification Mapping}

Table~\ref{tab:mapping} provides the exhaustive mapping from \texttt{PolicyCode} to
transmission channel, organized by \texttt{PolicyCategory}.

\begin{small}
\begin{longtable}{p{0.6cm} p{5.8cm} p{1.2cm} p{5.5cm}}
\caption{Classification of PolicyCodes to Transmission Channels}
\label{tab:mapping} \\
\toprule
\textbf{Code} & \textbf{Description} & \textbf{Channel} & \textbf{Rationale} \\
\midrule
\endfirsthead
\toprule
\textbf{Code} & \textbf{Description} & \textbf{Channel} & \textbf{Rationale} \\
\midrule
\endhead
\midrule
\multicolumn{4}{r}{\textit{Continued on next page}} \\
\bottomrule
\endfoot
\bottomrule
\endlastfoot

\multicolumn{4}{l}{\textit{Category 1: Revenue Measures to Protect Businesses}} \\
\midrule
1  & Accelerated asset depreciation (CIT) & CP & Reduces firm tax burden; preserves cash flow for operations \\
2  & Extend loss carry-forward (CIT) & CP & Allows firms to offset losses against prior profits; preserves solvency \\
3  & Broaden tax deductibility & CP & Reduces effective tax rate on firms; preserves operating margin \\
4  & Tax credits (business) & CP & Reduces firm tax liability; preserves cash flow \\
5  & Deferral of tax filing (business) & CP & Liquidity preservation (Decision Rule~3) \\
6  & Deferral of tax payments (business) & CP & Liquidity preservation (Decision Rule~3) \\
7  & Tax rate reduction (business) & CP & Reduces structural cost burden on firms; preserves firm viability \\
8  & Tax amnesty / incentives (business) & CP & Removes penalty-driven insolvency risk \\
9  & Accelerating refunds (business) & CP & Injects liquidity to firms from already-owed government obligations \\
10 & Lower advance payment & CP & Reduces prepayment burden; liquidity preservation \\
11 & Suspend debt collection & CP & Prevents enforcement-driven firm failure \\
12 & Suspend audit activities & CP & Administrative relief; reduces compliance burden during crisis \\
13 & Tax exemption (business) & CP & Eliminates specific tax obligations; preserves firm cash flow \\
14 & Tax waiver / suspension (business) & CP & Includes bankruptcy moratoriums; prevents firm exit \\
\midrule

\multicolumn{4}{l}{\textit{Category 2: Revenue Measures to Protect Individuals}} \\
\midrule
15 & Deferral of tax filing (individual) & CP & Liquidity preservation (Decision Rule~3) \\
16 & Deferral of tax payments (individual) & CP & Liquidity preservation (Decision Rule~3) \\
17 & Tax rate reduction (individual) & DI & Permanent increase in disposable income; standard multiplier channel \\
18 & Tax amnesty / incentives (individual) & DI & Increases disposable income of households \\
19 & Broaden tax deductibility (individual) & DI & Increases disposable income via lower effective tax rate \\
20 & Tax credits (individual) & DI & Direct increase in household disposable income \\
21 & Tax exemption (individual) & DI & Eliminates tax obligation; raises disposable income \\
22 & Tax waiver / suspension (individual) & DI & Reduces household tax burden; increases disposable income \\
\midrule

\multicolumn{4}{l}{\textit{Category 3: Revenue Measures to Boost Consumption/Demand}} \\
\midrule
25 & Exemption of consumption tax rates & DI & Reduces consumer prices; stimulates aggregate expenditure \\
26 & Lower consumption tax rates (VAT) & DI & Reduces consumer prices; transmits through spending channel \\
\midrule

\multicolumn{4}{l}{\textit{Category 4: Expenditure Measures to Boost Consumption/Demand}} \\
\midrule
27 & Investment into infrastructure & DI & Government spending multiplier; direct demand injection \\
28 & Green investments & DI & Government spending multiplier; direct demand injection \\
29 & Reactivate local tourism & DI & Consumption stimulus targeting specific sectors \\
\midrule

\multicolumn{4}{l}{\textit{Category 5: Health Expenditure Measures}} \\
\midrule
30 & Lower tax rates for medical items & H & Health system cost; calibrates $c_H$ \\
31 & Supply of low-cost medical items & H & Health system cost; calibrates $c_H$ \\
32 & Supply of high-cost medical items & H & Health system cost; calibrates $c_H$ \\
33 & Targeted health infrastructure & H & Health system cost; calibrates $c_H$ \\
34 & Expansion of human resources (health) & H & Health system cost; calibrates $c_H$ \\
gen. & General health related spending & H & Health system cost; calibrates $c_H$ \\
\midrule

\multicolumn{4}{l}{\textit{Category 6: Expenditure Measures to Individuals}} \\
\midrule
35 & Direct cash transfers for individuals & DI & Canonical demand injection; direct income flow to households \\
36 & Expansion of unemployment benefits & DI & Increases disposable income of unemployed; multiplier channel \\
37 & Temporary expansion of existing benefits & DI & Augments existing transfer programs; income flow to households \\
38 & Supplementary ad hoc programs & DI & Pandemic-specific income support to individuals \\
39 & Wage compensation / enhanced paid leave & CP & Dominant mechanism: preserves employer--employee matches (Kurzarbeit, CJRS, JobKeeper, Activit\'{e} Partielle, NOW). See Section~\ref{sec:code39}. \\
\midrule

\multicolumn{4}{l}{\textit{Category 7: Credit and Equity Measures}} \\
\midrule
40 & Preferential loans to firms & CP & Preserves firm solvency through liquidity provision \\
41 & Equity or quasi-equity investments & CP & Preserves firm capital structure; prevents insolvency \\
42 & Preferential loans to households & DI & Provides liquidity to households; transmits through spending \\
43 & Guarantees for loans for business & CP & Preserves firm access to credit; prevents insolvency cascade \\
\midrule

\multicolumn{4}{l}{\textit{Category 8: Expenditure Measures for Businesses}} \\
\midrule
45 & Income support general (business) & CP & Preserves firm operations; prevents exit of viable businesses \\
46 & Income support non-profit & CP & Preserves organizational capacity of non-profit sector \\
47 & Income support specific (sectoral) & CP & Preserves sector-specific productive capacity \\
48 & Income support education sector & CP & Preserves institutional capacity in education \\
\midrule

\multicolumn{4}{l}{\textit{Category 9: Transfers to Local Governments}} \\
\midrule
49 & Income support for local governments & CP & Preserves public service delivery capacity; prevents layoffs \\

\end{longtable}
\end{small}


\section{Treatment of Special Cases}

\subsection{Code 39: Wage Compensation / Enhanced Paid Leave}
\label{sec:code39}

Code~39 is classified as CP at the code level. The empirical justification is as
follows. Among the 140 measures coded as 39 across 37 countries, the overwhelming
majority ($>$90\% by fiscal volume) are short-time work schemes (Kurzarbeit, CJRS,
Activit\'{e} Partielle, NOW, JobKeeper) and apprenticeship wage subsidies---instruments
whose defining feature is that the payment flows \emph{through} the employer to
maintain the employment relationship. A small number of entries (e.g., parental
income compensation for school closures in Germany, pandemic leave payments in
Australia) are structurally closer to DI, as they compensate individual income loss
without preserving a productive firm--worker match.

Code-level classification is chosen over case-by-case assignment for three reasons.
First, it ensures reproducibility: any researcher applying the protocol to the same
dataset will obtain identical results. Second, the misclassification bias is bounded:
the DI-like entries within Code~39 represent less than 10\% of the code's fiscal
volume and are quantitatively small relative to the dominant short-time work programs.
Third, the sensitivity of results to this decision can be tested by reclassifying
the flagged entries (see Section~\ref{sec:flags}).

\subsection{Code 45: Income Support General (Business)}

Code~45 (324 observations, 0.64\% of GDP) contains predominantly firm-level grants
and income support: German Soforthilfe, French Fonds de Solidarit\'{e}, and similar
programs across OECD economies. A small number of US entries (notably CARES Act
State \& Local Funding) involve intergovernmental transfers with mixed purpose. These
are classified as CP at the code level, consistent with the code's definition
(``income support general'' under Category~8: Expenditure Measures for Businesses).

\subsection{Tax Deferrals (IMF Category 3)}

All 154 tax deferral measures are classified as CP regardless of the recipient
(business or individual). This reflects the economic nature of deferrals as
temporary liquidity injections: the government provides an implicit zero-interest
loan that must be repaid. The transmission mechanism is capacity preservation
(preventing liquidity-driven default), not demand injection (no permanent increase
in disposable income). The fiscal cost of deferrals is also structurally different
from both DI and CP: the budgetary impact is limited to the time value of the
deferred revenue and any non-repayment, not the face value of the deferral. This
is reflected in the IMF's below-the-line treatment.

Note that this classification decision implies that the effective $\kappa_F^{CP}$
for the deferral component of CP is substantially below the weighted average
$\kappa_F^{CP}$, since the realized fiscal cost of deferrals approaches zero for
solvent taxpayers. Sensitivity analysis should examine results under exclusion of
deferrals from CP.

\subsection{Health Expenditures (Category H)}

Health measures (Codes 30--34 and \texttt{general}; 235 observations, $\sim$0.43\% of
GDP) are excluded from both DI and CP. In the model, these expenditures are not a
choice variable in the planner's fiscal composition problem but an endogenous cost
$c_H \theta_k$ driven by the epidemiological state. The classified H measures
provide the empirical basis for calibrating~$c_H$.

\subsection{Codes 23 and 24}

Code~23 (1 observation: tax-free alcohol for disinfectant production, Germany) and
Code~24 (5 observations: customs duty relief on medical goods) are health-adjacent
measures. They are classified as H, consistent with their purpose of facilitating
medical supply availability during the pandemic.

\subsection{Policy Extensions (\texttt{broad\_fiscal} = 2)}

Policy extensions (443 observations) carry \texttt{broad\_fiscal\_size} = 0 and
document the temporal extension of existing programs without additional fiscal
commitment. Each extension inherits the transmission channel of its base
\texttt{PolicyCode}. Extensions are flagged (\texttt{is\_extension} = 1) so they can
be excluded from volume calculations while preserving the timeline information.

\subsection{EU-Level Measures (\texttt{broad\_fiscal} = 3)}

The 29 NextGenerationEU/RRF entries are classified based on their
\texttt{PolicyCode} (all coded as 32, Category~5) but carry a separate flag
(\texttt{is\_eu} = 1). These supranational measures differ from national fiscal
responses in timing (disbursement lagged), financing (EU-level debt), and
conditionality (milestone-based). They should be analyzed separately or excluded
from the baseline specification.

\subsection{Code \texttt{other} (41 Observations)}
\label{sec:other}

The residual category \texttt{other} (41 observations, all policy extensions with
zero fiscal weight) is classified case-by-case based on the measure description.
Since all entries have \texttt{broad\_fiscal\_gdp} = 0, misclassification has no
quantitative effect on the empirical analysis. The R implementation script assigns
these based on keyword matching and manual verification of ambiguous cases. A
complete listing of the case-by-case assignments is provided in the R script
documentation.


\section{Robustness Flags}
\label{sec:flags}

The classification adds the following binary flags for sensitivity analysis:

\begin{itemize}[leftmargin=2em]
  \item \texttt{is\_extension}: 1 if \texttt{broad\_fiscal} = 2 (zero fiscal weight).
  \item \texttt{is\_eu}: 1 if \texttt{broad\_fiscal} = 3 (supranational measure).
  \item \texttt{is\_deferral}: 1 if IMF \texttt{category} = 3 (tax deferral).
  \item \texttt{is\_blw}: 1 if IMF \texttt{category} = 2 (below-the-line: loans and
    guarantees).
\end{itemize}

\noindent These flags enable the following robustness checks:

\begin{enumerate}[leftmargin=2em]
  \item \textbf{Exclude deferrals from CP}: reclassify \texttt{is\_deferral} = 1
    measures as a separate category or exclude entirely, testing whether the CP
    effectiveness parameter $\eta$ is sensitive to the inclusion of pure timing
    measures.
  \item \textbf{Exclude extensions}: restrict analysis to
    \texttt{broad\_fiscal} = 1 (core fiscal measures with positive fiscal weight).
  \item \textbf{Exclude EU measures}: restrict to national fiscal responses.
  \item \textbf{Separate above- and below-the-line CP}: partition CP into direct
    measures (grants, subsidies; \texttt{is\_blw} = 0) and contingent measures
    (loans, guarantees; \texttt{is\_blw} = 1), testing whether $\eta$ and
    $\kappa_F^{CP}$ differ across sub-categories.
\end{enumerate}


\section{Validation}

The classification can be cross-validated against two external benchmarks:

\begin{enumerate}[leftmargin=2em]
  \item \textbf{IMF fiscal accounting}: Above-the-line measures
    (\texttt{category} = 1) should appear predominantly in both DI and CP;
    below-the-line measures (\texttt{category} = 2) should be exclusively CP
    (loans and guarantees preserve firm capacity). Tax deferrals
    (\texttt{category} = 3) should be exclusively CP per the protocol.
  \item \textbf{Aggregate consistency}: The country-level DI and CP aggregates
    (as \% of GDP) should be broadly consistent with the IMF's decomposition of
    above-the-line versus below-the-line spending, adjusted for the
    transmission-based reclassification of certain above-the-line measures.
\end{enumerate}





\end{document}
